\chapter{The Model: Open Weights as the Commodity Layer}
\label{ch:model}

The first element of the AI war is the model itself. The Western frontier is closed: the best US models are accessible only as metered APIs, their weights corporate secrets, their pricing a rent on capability. The Chinese frontier is open: weights downloadable, licenses mostly permissive (MIT, Apache~2.0)---though, as \S\ref{sec:model-survey} details, the newest flagships have begun attaching custom commercial terms that are themselves evidence for this book's rent-relocation thesis---and API prices that until recently sat at or near marginal inference cost. That last proposition is the one this chapter has to qualify hardest, because the price sheets of July and August 2026 stopped obeying it: the American incumbent's cheapest tier now undercuts most of the Chinese frontier on list price, while the Chinese price leader has \emph{raised} prices by as much as elevenfold. Section~\ref{sec:price-inversion} works through what that inversion does and does not mean. This chapter surveys the whole picture---the labs, the releases, the algorithmic innovations, the quantization ecosystem that is collapsing the last barrier between frontier capability and commodity hardware, the price line at which the two ecosystems now meet, and---new at this currency date---the licence clause by which the makers of the commons have begun to negotiate a share of the American clouds' revenue from serving it (Section~\ref{sec:toll-road}).

\section{The DeepSeek shock and the price war (2024--2025)}

The opening move replicated the solar price collapse, compressed into eighteen months. DeepSeek---a 2023 spin-out of the High-Flyer quantitative fund---released V2 in May 2024 with Multi-head Latent Attention and a fine-grained MoE at roughly RMB~1--2 per million tokens; within two weeks ByteDance, Zhipu, Alibaba, and Baidu had cut prices by up to 97\%, Baidu making its mainline models free within four hours of Alibaba's cut~\cite{price-war}. V3 (December 2024)---671B parameters, 37B active, and the first open large-scale model natively trained in FP8---reached GPT-4-class capability with a published marginal training cost of \$5.6 million (2.79 million H800-hours), a figure legitimately contested as excluding R\&D and infrastructure, but directionally transformative either way~\cite{deepseek-v3}. R1 (January 2025) matched OpenAI's o1-class reasoning, was released under MIT license, topped the US iOS App Store, and erased \$589 billion of NVIDIA's market capitalization in a single trading day---the largest one-day single-company loss in US market history~\cite{r1-shock}. The peer-reviewed Nature version disclosed that R1's reinforcement-learning stage cost about \$294,000~\cite{r1-nature}.

The strategic content of the shock was mispriced by markets and correctly priced by Beijing: it demonstrated that frontier capability could be commoditized, that export-control-constrained hardware forced \emph{algorithmic} efficiency (MLA, extreme sparsity, FP8 training, GRPO, multi-token prediction) which then compounds on any hardware, and that open weights convert capability into global adoption instantly.

A cost-accounting discipline applies to every training-cost figure in this book. A ``\$5.6M training run'' is final-run marginal compute [P]; the full program cost adds failed and ablation runs, post-training and RL, synthetic-data inference, data acquisition and cleaning, personnel, hardware depreciation, networking, power, and the prior research that produced the reusable architecture. DeepSeek's own disclosed hardware position ($\sim$50,000 H800s; buildout estimates to \$1.6B [A]) bounds the true denominator~\cite{deepseek-v3}. The strategic point survives the correction---marginal-cost transparency itself resets buyer expectations---but the book does not treat run cost as program cost.

\section{An innovation genealogy}
\label{sec:genealogy}

A release chronology shows cadence but not authorship, and the strategically important question is which architectural innovations are Chinese-origin, which were independently reproduced, and---the question this book cares about most---which \emph{reduce a hardware requirement}. Table~\ref{tab:genealogy} attempts the genealogy.

\begin{table}[htbp]
  \centering
  \scriptsize
  \caption{Innovation genealogy of the open frontier. ``Hardware effect'' names the resource the technique economizes---the column that connects this chapter to Chapter~\ref{ch:compute}.}
  \label{tab:genealogy}
  \begin{tabular}{p{31mm}p{20mm}p{34mm}p{48mm}}
    \toprule
    Technique & First at scale & Origin / diffusion & Hardware effect \\
    \midrule
    Multi-head latent attention (MLA) & DeepSeek V2, 2024-05 & Chinese-origin; widely adopted & KV cache to $\sim$2\% of dense: \emph{memory capacity and bandwidth} \\
    \addlinespace
    Fine-grained + shared-expert MoE & DeepSeek V2/V3 & Chinese-origin refinement of a Western idea & activated fraction to 1.8--3\%: \emph{FLOPs and bandwidth} \\
    \addlinespace
    Native FP8 pretraining & DeepSeek V3, 2024-12 & first open large-scale demonstration & \emph{memory and compute}; sets the domestic-chip numeric format \\
    \addlinespace
    GRPO and critic-free RL & DeepSeek R1, 2025-01 & Chinese-origin; broadly reproduced & removes critic model: \emph{post-training compute} \\
    \addlinespace
    Multi-token prediction & DeepSeek V3 & Chinese-origin at scale & \emph{tokens per step}; speculative decode \\
    \addlinespace
    Sparse attention (DSA) & DeepSeek V3.2, 2025-09 & Chinese-origin & $O(L^2)\!\to\!O(Lk)$: \emph{long-context compute}---the decisive lever for CPU-class prefill \\
    \addlinespace
    INT4 / MXFP4 QAT release formats & Kimi K2 Thinking 2025-11; K3 2026-07 & Chinese-origin at frontier scale; gpt-oss parallel & \emph{serving memory and bandwidth}; 4-bit as release format \\
    \addlinespace
    Hybrid linear attention at frontier scale & Kimi K3, 2026-07; GLM-5.3-Flash, 2026-08 & Chinese-origin & \emph{long-context memory traffic}; Zhipu claims $\sim$4.4$\times$ KV-cache reduction against its own flagship [P, vendor] \\
    \addlinespace
    Expert-parallel dispatch/communication overlap & DeepSeek V3 infra & Chinese-origin engineering & \emph{interconnect}: tolerates weaker fabric \\
    \bottomrule
  \end{tabular}
\end{table}

Read down the right-hand column, but read it carefully, because it is easy to overstate what this table shows, in both its arithmetic and its causality. The arithmetic first: it is \emph{not} true that almost none of these techniques economizes dense FLOPs---the table itself refutes it. Fine-grained MoE reduces active FLOPs directly; FP8 reduces compute cost as well as memory; critic-free RL removes an entire model's worth of post-training compute; multi-token prediction and speculative decoding change compute utilization per emitted token. The accurate reading is one of \emph{emphasis}, not exclusion: the portfolio is concentrated unusually strongly on memory capacity, memory bandwidth, interconnect efficiency, and low-precision execution---the resources that became disproportionately scarce under export control---while the compute savings it also delivers are the kind any laboratory, constrained or not, would pursue.

The causality second. Efficiency is valuable to everyone; MoE, low precision, sparse attention, and communication overlap were all being advanced by unconstrained Western laboratories in the same period, and most would likely have progressed without any export controls at all. The defensible claim is therefore \emph{constraint-consistency}, not unique causation: export controls are a plausible accelerant and a selection pressure that shaped which efficiencies were prioritized and how aggressively, but not the sole origin of the techniques. Table~\ref{tab:causality} states the questions that separate strong from weak co-design evidence, and applies them.

\begin{table}[htbp]
  \centering
  \scriptsize
  \caption{A causality matrix for the genealogy: what would make ``co-design under sanction'' strong or weak for a given technique, with the two clearest verdicts from Table~\ref{tab:genealogy}.}
  \label{tab:causality}
  \begin{tabular}{p{56mm}p{42mm}p{34mm}}
    \toprule
    Diagnostic question & Interpretive value & Verdict on examples \\
    \midrule
    Did the technique predate the relevant restriction? & Weakens a sanctions-origin claim & MoE, low precision: yes---weak \\
    \addlinespace
    Was it independently pursued by unconstrained labs? & Shows general value, not uniquely Chinese adaptation & MoE, FP8, sparse attention: yes---weak \\
    \addlinespace
    Does it specifically match a domestic hardware limitation? & Strengthens co-design attribution & MLA, dispatch/comm overlap: yes---strong \\
    \addlinespace
    Is there an explicit domestic numeric-format, compiler, or processor target? & Direct evidence of hardware--algorithm co-design & UE8M0 FP8 ``for next-generation domestic chips'': yes---strongest \\
    \addlinespace
    Would it remain commercially attractive without sanctions? & Separates national adaptation from ordinary cost reduction & nearly all: yes---caps every claim above \\
    \bottomrule
  \end{tabular}
\end{table}

On this matrix, UE8M0 FP8 with a declared domestic-silicon target is strong co-design evidence; generic MoE adoption is weak; MLA and the interconnect-tolerance engineering sit in between, matching domestic limitations specifically enough to carry weight. The genealogy's strategic conclusion, stated at the strength the matrix actually supports: the \emph{direction} of Chinese architectural emphasis is exactly the direction that makes the hardware of Chapter~\ref{ch:compute} viable, and export controls demonstrably sharpened that emphasis even where they did not originate it.

\section{The ecosystem survey: eight frontiers, one commons}
\label{sec:model-survey}

Consistent with playbook step~4 (selection, not planning; move~M4), the state did not anoint DeepSeek; it flooded the sector. Table~\ref{tab:model-timeline} compresses the release cadence of the major laboratories; the notes that follow give each lab's trajectory. Two structural observations frame the table. First, the cadence itself is the weapon: the interval between a US closed-frontier release and its Chinese open-weight counterpart compressed from six--eight months in 2024 to ``a matter of weeks'' by mid-2026~\cite{aei-2026}. Second---stated in its properly qualified form---open-weight release has become the dominant international distribution strategy of the leading Chinese laboratories, \emph{although license terms range from permissive Apache/MIT grants to custom commercial licenses that preserve control over high-revenue deployment}. The table therefore carries a license column, because the license is strategic data, not legal boilerplate: Kimi K3, the largest open-weight model ever released, ships under a custom license that permits use, modification, redistribution, and derivatives but requires a separate commercial agreement from qualifying model-as-a-service businesses above a revenue threshold, with attribution conditions on very large deployments~\cite{kimiLicense}. That is this book's own industrial-organization thesis enacted in a license file: opening the weights does not abolish rent---it relocates it to large-scale commercial service, where the releasing lab retains a claim. By \datafreeze the pattern has three frontier-scale instances with three different thresholds and two different instruments, the newest of which substitutes a security review for a fee (\S\ref{sec:glm53}).

Two terminological disciplines apply throughout. These are \emph{open-weight} releases, not open-source: with few exceptions the training code, data, and full recipes are not published, which is why R1 spawned an entire replication literature~\cite{r1-replication}. And the correct narrow claim is that leading Chinese laboratories have made open-weight release the dominant \emph{international distribution strategy} for their frontier-adjacent families---not that every Chinese frontier system is open (several are not), nor that the West releases nothing (gpt-oss, Gemma, OLMo, and NVIDIA's Nemotron line are real, if outclassed at the top end).

\begin{table}[htbp]
  \centering
  \scriptsize
  \caption{Chinese open-weight frontier releases, 2024--August 2026 (selected; parameters as total/active where MoE; benchmark figures are vendor-reported [P] unless noted). License classes: Apache = Apache 2.0; MIT; custom = custom commercial license (weights downloadable, high-revenue service deployment separately controlled). ``AA index'' is the Artificial Analysis Intelligence Index, always with its version: v4.1.1 is the launch-week board of August 2026, v4.3 the re-based board of 7~September~\cite{aa-v43}; scores are not comparable across versions. Current as of \datafreeze (the edition date); the most recent cycle is analyzed in Appendix~\ref{app:bulletin}.}
  \label{tab:model-timeline}
  \begin{tabular}{p{15mm}p{17mm}p{28mm}p{14mm}p{58mm}}
    \toprule
    Date & Lab & Model & License & Significance \\
    \midrule
    2024-05 & DeepSeek & V2 (236B/21B) & MIT-class & MLA + fine-grained MoE; triggers national LLM price war \\
    2024-12 & DeepSeek & V3 (671B/37B) & MIT-class & native FP8 pretraining; \$5.6M marginal cost claim \\
    2025-01 & DeepSeek & R1 (671B/37B) & MIT & o1-class reasoning; the NVIDIA crash \\
    2025-04 & Alibaba & Qwen3 (0.6B--235B/22B) & Apache & family of eight, 36T tokens, hybrid reasoning \\
    2025-06 & Baidu & ERNIE 4.5 (to 424B) & Apache & wholesale open-sourcing of a former closed flagship \\
    2025-07 & Moonshot & Kimi K2 (1T/32B) & modified MIT & largest open model to date; agentic focus \\
    2025-08 & DeepSeek & V3.1 & MIT & hybrid think/non-think; UE8M0 FP8 ``for next-generation domestic chips'' \\
    2025-09 & DeepSeek & V3.2 (DSA) & MIT & sparse attention: $O(L\cdot k)$; API cut $>$50\% again \\
    2025-11 & Moonshot & K2 Thinking (1T/32B) & modified MIT & native INT4 QAT; beats GPT-5 on HLE (44.9 vs 41.7) \\
    2026-02 & Zhipu & GLM-5 (744B/40B) & MIT & SWE-bench Verified 77.8\%; Ascend-only training \emph{reported} in the press (Jun.\ 2026), no primary located [A/C] \\
    2026-02 & MiniMax & M2.5 (230B/10B) & custom & \#1 Multi-SWE-Bench; \$0.15/\$1.20 per Mtok \\
    2026-04 & DeepSeek & V4 (1.6T/49B) & MIT & 1M context; $\sim$150$\times$ cheaper than GPT-5.5 output \\
    2026-04 & Alibaba & Qwen3.6 (27B dense; 35B-A3B) & Apache & Qwen crosses 1B cumulative downloads \\
    2026-04 & Xiaomi & MiMo-V2.5 (310B/15B); V2.5-Pro (1.02T/42B) & MIT & omni-modal input, 1M context; a handset maker ships a trillion-parameter open model~\cite{mimo-v25} \\
    2026-06 & MiniMax & M3 ($\sim$428B/23B) & custom & multimodal; 1M context; AA index 44 at release (June board, pre-re-basing) \\
    2026-06 & Zhipu & GLM-5.2 ($\sim$753B/40B) & MIT & \#1 open model on the June Artificial Analysis board \\
    2026-07 & Moonshot & Kimi K3 (2.8T/$\sim$104B) & custom & largest open model ever; MXFP4; global top-5 overall \\
    2026-07 & DeepSeek & V4-Flash ($\sim$280B/13B) & MIT & weights shipped; 1M context; \$0.14/\$0.28 per Mtok; gains from post-training alone \\
    2026-08 & Alibaba & Qwen3.8-Max (2.4T/$\sim$95B) & custom & \textbf{weights shipped}; \$50M revenue-threshold licence; text-only; 262K native context (\S\ref{sec:maxtier}) \\
    2026-08 & Alibaba & Qwen3.8-27B & Apache & vision-language companion; the permissive half of the release \\
    2026-08 & Zhipu & GLM-5.3 ($\sim$744B/40B) & custom & weights released 28 Aug.; \texttt{glm-5.3} licence, reported security review above \$10B MaaS revenue; AA index 45 on v4.3 (60, \#2 of 111, at launch on v4.1.1) (\S\ref{sec:glm53}) \\
    2026-08 & Zhipu & GLM-5.3-Flash (320B/18B) & MIT & the stealth model ``Ox~Alpha''; natively multimodal, 1M context; AA index 42 on v4.3 at \$0.25 per task (57, \#4, at launch on v4.1.1); \emph{served} on undisclosed domestic silicon (\S\ref{sec:glm53}) \\
    2026-08 & Meta & Muse Glimmer (30B dense) & Apache & Meta's return to open weights after 16 months; open Spark 1.2 promised (\S\ref{sec:muse}) \\
    2026-08 & Xiaohongshu & dots3-note Preview (280B/16B) & Apache & text/image/video/audio input, 512K context; TEMPO long-horizon RL for agents; a social platform ships an agent model~\cite{dots3-note} \\
    \bottomrule
  \end{tabular}
\end{table}

Two disciplines govern this table. \emph{First, a currency date}: every fact in the stable chapters is current as of \datafreeze---the date of this edition, printed on its title page---and nothing learned after that date is retrofitted into the analysis. Freezing the data at a date \emph{earlier} than the edition invites exactly the contradiction a reader would expect---a declared July cutoff sitting a few pages from an August release---so the rule adopted here is the only one that cannot contradict itself: \emph{the cutoff is the edition date}. Correspondingly, Appendix~\ref{app:bulletin} carries the most recent cycle in detail so that the fastest-moving material sits in one replaceable place rather than being threaded through the argument. Applying the rule, the table above runs through the release of GLM-5.3's weights on 28~August and the reveal, the day before, that the stealth model ``Ox~Alpha'' was GLM-5.3-Flash~\cite{glm53,glm53-flash}; Qwen3.8-Max sits at its 3~August announcement~\cite{qwen-announcement,qwen38-max}, and DeepSeek's V4-Flash where its actual 31~July release date puts it~\cite{v4flash}. No analytical claim in this book depends on which model holds the crown in any given week; the three-artifact structure---stable monograph, machine-readable claim register, living appendix---exists precisely so that the argument does not chase the cadence it describes. \emph{Second, exact naming}: the open Qwen3.6 releases are the 27B dense and 35B-A3B variants under Apache 2.0~\cite{qwen36-repo}---``Qwen3.6'' is not one frontier-scale flagship, and the table says so.

\emph{Second, the breadth of the table is itself the datum.} The rows are not the work of one national champion and a few followers. DeepSeek, Alibaba, Moonshot, Zhipu, MiniMax, Baidu, Tencent and Huawei are expected entrants; the last two rows added in 2026 are not. Xiaomi---a handset and appliance maker---open-sourced a 310B/15B omni-modal model and a 1.02-trillion-parameter text model under MIT in April~\cite{mimo-v25} [V]. Xiaohongshu---a lifestyle and commerce platform of more than 400 million monthly users, preparing a Hong Kong listing at a reported valuation above \$70B---released \emph{dots3-note Preview} on 14~August under Apache~2.0: 280B total and 16B active parameters, text, image, video and audio input, a 512K context, and a reinforcement-learning method (TEMPO, Chapter~\ref{ch:harness}) built for agent runs of ten hours or more; the same laboratory's dots-note-3.0 had scored a perfect 42/42 at the 2026 International Mathematical Olympiad three weeks earlier, and Huawei's Ascend stack carried the open model on day zero~\cite{dots3-note,dots-imo,dots3-ascend,rednote-ipo} [V/A]. What that pair of rows records is that frontier-scale model development in China has become a capability a large technology platform is \emph{expected} to have, rather than the specialty of a handful of laboratories---which is the industrial-base observation of Part~I arriving at the model layer, and a stronger fact about the ecosystem than any single release in the table.

\subsection{The open-weight ratchet: why the commons is now an equilibrium}
\label{sec:maxtier}

The Qwen3.8-Max release deserves more than a chronology entry, because it converts this book's central claim from a description of state strategy into something considerably stronger: a description of a \emph{competitive equilibrium that no participant can profitably defect from}.

Start with the fact pattern. Across the Qwen family's history the numbered open releases (7B, 72B, the Qwen3 family, the 3.6 dense and A3B variants) shipped under Apache~2.0 and drove the download statistics that make Qwen the most-adopted open family on earth---roughly a billion cumulative downloads by March 2026, 153.6 million in February alone, more than half of all global open-model downloads, and over 200,000 derivative variants~\cite{qwen-downloads}. The \emph{Max} tier, meanwhile---Qwen2.5-Max, Qwen3.5-Omni, Qwen3.6-Plus, Qwen3.7-Max---had been consistently proprietary: API and chatbot only. That was the settled pattern, and a reasonable observer in June 2026 would have predicted it continuing.

It did not continue. On 27 July, Moonshot open-sourced Kimi K3's weights at 2.8 trillion parameters. Seven days later Alibaba announced Qwen3.8-Max and committed, publicly and unambiguously, on its own official account: \emph{``Next week, the open weights of Qwen3.8-Max will be released, and Qwen3.8-27B is also going open-weights''}~\cite{qwen-announcement}. The first Max-tier model in the family's history to carry an open-weight commitment carried it within a week of a rival opening at comparable scale---alongside a published blog with the full benchmark table, and a complete project trace on GitHub for the flagship autonomy demonstration~\cite{qwen38-blog}.

The causal reading is hard to avoid, and it is the most important structural observation in this chapter. \emph{Once one Chinese laboratory opens at frontier scale, the others cannot stay closed.} The mechanism is not ideology and not state compulsion; it is the arithmetic of the strategy every Chinese lab is running. Their competitive position is built on adoption---downloads, derivatives, developer mindshare, standards gravity, and the deployment volume that Chapter~\ref{ch:theory}'s flywheel converts into learning and cost advantage. Against a rival whose frontier-scale weights are downloadable, a metered-only flagship forfeits precisely the currency the strategy is denominated in. Whatever premium the closed tier collects is smaller than the adoption it surrenders, because the adoption is what the whole apparatus---hardware demand, standards position, national ecosystem---was built to capture. Closing your best model, in that ecosystem, is unilateral disarmament.

This is a \emph{ratchet}, and its properties matter. Each open release at a new capability level raises the floor every competitor must match; the floor never falls, because no lab gains by being the one that retreats behind an API while its rivals ship weights. Defection is individually irrational even when it would be collectively profitable---the exact inverse of the Western equilibrium, where each frontier laboratory's rivalry is over premium capability sold at metered access, and closing is individually rational for each. Two ecosystems, two stable equilibria, each internally coherent, each self-enforcing without a planner. The Chinese one produces a commons; the American one produces a rent.

Three consequences follow. \emph{First}, the thesis is more robust than a policy-grounded account of openness would be. Grounding Chinese openness in the Five-Year Plan, MIIT guidance, and WAICO---the natural reading, and the common one---invites the obvious rebuttal that policy can reverse. The ratchet does not depend on Beijing's continued preference. Even if the state's enthusiasm cooled tomorrow, a Chinese lab that closed its frontier tier would hand the developer base, the derivative ecosystem, and the standards position to whichever rival stayed open. \emph{Second}, it predicts the pattern forward: the next frontier-scale Chinese release will be open, and the one after that, not because anyone decreed it but because the first defector loses---a prediction the next release tested within the month, when Zhipu published GLM-5.3's weights on 28~August under a custom licence and its distilled companion under MIT (\S\ref{sec:glm53}). \emph{Third}, it refines rather than contradicts the rent-relocation argument of Section~\ref{sec:io}: rent still relocates---to licence terms (Kimi K3's revenue threshold), to hosted service, to hardware, to the complements of Table~\ref{tab:rent}---but it relocates \emph{out of the weights}, and the ratchet is why it cannot relocate back.

The register carries the near-term test, and it has now resolved---in two directions at once, which is why it was registered as two forecasts rather than one. \emph{The weights shipped.} Within roughly ten days of the announcement, Alibaba published the 2.4-trillion-parameter model to Hugging Face---as \texttt{Qwen3.8-2.4T-A95B}, not under the product name ``Max''---at 4.89\,TB across 213 BF16 shards, with an FP8 variant, followed on 15 August by \texttt{Qwen3.8-27B}~\cite{qwen38-weights,qwen38-scmp}. The 90\% forecast resolves \textbf{YES}. \emph{The licence is not permissive.} The flagship carries a custom \texttt{qwen3.8-max} licence requiring a separately negotiated commercial agreement from any licensee whose model-as-a-service or AI-work-assistant business exceeds \textbf{US\$50 million of revenue in any consecutive twelve months}, with an additional attribution obligation above 100 million monthly active users or US\$20 million of monthly revenue; internal use whose outputs are not shared with third parties is unrestricted~\cite{qwen38-weights}. Only the 27B companion is Apache~2.0. The 55\%$\to$30\% permissive-licence forecast therefore resolves \textbf{NO} on the model that matters, and the correction is worth printing loudly because much of the trade press reported the release as uniformly Apache-licensed and was wrong~\cite{qwen38-scmp}.

Two further precisions the record requires. The Max-tier model is \emph{text-only}, not multimodal as the pre-release press consensus had it; the 27B is the vision-language variant. And the licence contains \emph{no geographic restriction}---the US, EU, UK and South Korea exclusion that circulated in early August belongs to MiniMax's H3 video model, a different company and a different artifact, and conflating them would misstate both.

So the ratchet holds on weights and fails on terms, which is precisely the discrimination the two forecasts were built to make. The mechanism compels \emph{weights onto the commons}, because adoption is the currency; it does not compel Apache terms, because rent relocation operates on exactly that margin. What Alibaba has built is a commons with a toll gate at commercial scale: free to researchers, to enterprises deploying internally, and to every derivative-builder below fifty million dollars---which is to say free to essentially the entire population that generates ecosystem gravity---and chargeable precisely where a rival could resell it as a service. Kimi K3 set the pattern with a twenty-million-dollar threshold; Alibaba raised the number and kept the structure. This is not openness retreating. It is openness that has learned where its own rent lives, and Section~\ref{sec:io}'s rent-relocation table predicted the destination; Section~\ref{sec:toll-road} follows the clause to the counterparty it has now found.

\subsection{The ratchet crosses the Pacific: Meta's return to open weights}
\label{sec:muse}

On 10 August 2026 the ratchet acquired its first American datapoint. Meta released the weights of \emph{Muse Glimmer}---a 30B-parameter dense multimodal model with a $\sim$128K context, distilled from its closed flagship and engineered to run on a single consumer GPU---under Apache~2.0, the most permissive licence Meta has ever attached to a frontier-family artifact, and a deliberate break from the custom community licences of the Llama era~\cite{muse-glimmer,muse-aa}. Zuckerberg announced the release personally on X and, in a 6{,}500-word letter the same day, recommitted Meta to ``American open source''---with a companion promise that an open-weight \emph{version} of the flagship Muse Spark 1.2 would follow ``in the coming weeks''~\cite{zuck-letter}. The sequence behind the reversal is worth restating precisely, because it is the argument of this chapter run in mirror image: Llama collapsed below 1\% of routed open-model volume; Meta shelved its Behemoth flagship, pivoted to closed development through late 2025, and shipped Muse Spark closed in April 2026; and then, sixteen months after its last open release---with Chinese families at roughly 61\% of routed tokens and the K3/V4-Flash/Qwen3.8 barrage dominating the commons---it returned to open weights~\cite{muse-glimmer,openrouter-2026}.

Does this accelerate the Chinese position or begin to reverse it? Both readings deserve their strongest form, and the evidence to date assigns them different time horizons.

\emph{The re-anchoring case.} An Apache-licensed American model at genuine quality, backed by Meta's distribution and compute, gives Western enterprises and governments what the trust chapter (Chapter~\ref{ch:trust}) argues they lack: a way to consume open weights without consuming Chinese provenance. If the promised open Spark 1.2 ships at competitive frontier quality, the commons stops being singularly Chinese---and with it the monoculture argument of Chapter~\ref{ch:order66} weakens, the certification problem eases for Western deployers, and the derivative ecosystem acquires a second pole. Zuckerberg's own framing is explicitly this: American open source must lead, and US policy should reduce the friction preventing it~\cite{zuck-letter}. For the non-aligned world of Chapter~\ref{ch:nonaligned}, a two-pole commons is strictly better than a one-pole commons, and this book has argued throughout that a contested commons is the healthiest available outcome.

\emph{The acceleration case.} Three facts cut the other way, and at present they cut deeper. First, \emph{what is actually open}: the small distilled model shipped; the flagship is a promise, of an unspecified ``version,'' under an unspecified licence---precisely the pattern this chapter documented for Alibaba's Max tier, but starting from sixteen months further behind in derivative ecosystems, fine-tune tooling, and deployed mindshare. Second, \emph{measured quality}: on the independent Artificial Analysis index as it stood in August (v4.1.1, before the September re-basings) Glimmer scored 35---above Gemma~4 31B at 30, below Qwen3.6 27B Reasoning at 38 and Kimi K2.5 at 36---with an 82\% measured hallucination rate against 49\% for the comparable Qwen~\cite{muse-aa}. A rule of method applies to that number and to every index number that follows: Artificial Analysis re-based the index twice in the week before \datafreeze, on 4~September (v4.2) and 7~September (v4.3), and the two re-basings moved every model's score by ten to fifteen points and re-ordered the top, so this book treats \emph{rank within a version} as informative and treats any comparison of scores across versions as meaningless, and prints the version beside every score~\cite{aa-v43}. The entry is competitive, not commanding; it raises the floor of the commons without capturing it. Third, and structurally most important, \emph{the ratchet interpretation}: Meta's return is what the mechanism of Section~\ref{sec:maxtier} predicts any lab whose strategy depends on adoption will eventually do. It is therefore \emph{evidence for the ratchet as a mechanism}---competitive compulsion, not policy preference, now drives openness on both sides of the Pacific---while leaving entirely open who owns the commons the ratchet enlarges. A commons contested by an American entrant validates the terrain Chinese strategy chose; it vindicates the fast-follower bet that value would migrate out of the weights; and every argument of Chapter~\ref{ch:theory} about rent relocation now applies to Meta too, which must monetize adjacent to weights it gives away, against incumbents seven releases into that same discipline.

The falsifiable statement, registered in Appendix~\ref{app:dashboard}: the re-anchoring case requires the open Spark 1.2 to ship within the committed window, at flagship-representative quality, under terms permissive enough to seed derivatives---three conditions, each observable within a quarter. If they hold, the openness scoreboard genuinely re-opens, and later editions of this book will say so. If the pattern instead repeats the Max-tier template---small model open, flagship deferred, terms restrictive---then August 2026 will read as the week the American ecosystem conceded that the Chinese strategy was correct and began running it from behind. Either way, one proposition was strengthened in August: \emph{nobody now argues the commons does not matter}. The contest has moved from whether to open to who opens more, faster, at higher quality---which is the contest the convergence strategy was built to win, on terrain where the incumbent advantage of distribution matters less than the compounding advantages of cadence and cost.

\subsection{GLM-5.3, ``Ox Alpha,'' and the serving claim}
\label{sec:glm53}

The last fortnight before the currency date supplied the sharpest available test of three propositions this chapter has been building---the ratchet, the deployable frontier, and the sovereignty of the stack underneath---and it resolved them in three different directions, which is why it is worth taking slowly.

\emph{The ratchet held on weights and moved again on terms.} Zhipu announced GLM-5.3 on 13--14 August, opened the API on 18 August, and published the weights on 28 August: $\sim$744B total parameters with $\sim$40B active, 78 layers, 256 routed experts with 8 active per token, a one-million-token context~\cite{glm53} [P]. The weights shipped, as \S\ref{sec:maxtier}'s mechanism predicts they must. The licence, however, is neither MIT nor Apache~2.0 but a custom \texttt{glm-5.3} instrument whose reported terms require a security review of model-as-a-service operators above \$10~billion of trailing revenue~\cite{glm53} [P/A on the reported threshold]. That is the third frontier-scale instance of the same structure and the most revealing of the three, because it changes the instrument as well as the number: Moonshot set a \$20~million revenue threshold and Alibaba a \$50~million one, both requiring a negotiated \emph{commercial} agreement, while Zhipu sets a threshold two orders of magnitude higher and asks not for money but for \emph{review}. Rent relocated to licence terms in the first two cases; in the third, what is retained at commercial scale is a governance hook. Either way the direction is the one Section~\ref{sec:io} predicted---out of the weights, into the terms of large-scale service---and the population below the threshold, which is to say every researcher, every enterprise deploying internally, and every derivative-builder on Earth, gets the artifact free.

\emph{The scoreboard moved, and the cheap model moved it further than the flagship.} On 26 August Zhipu released GLM-5.3-Flash, and on 27 August confirmed it had been running publicly as the unattributed stealth model ``Ox~Alpha'': 320B total parameters with 18B active, hybrid sparse and linear attention, the first natively multimodal member of the GLM-5 series, trained on a 30-trillion-token multimodal corpus, a one-million-token context---and released under \textbf{MIT}, at \$0.15/\$0.50 per million tokens, roughly $9\times$ below the flagship~\cite{glm53-flash} [P]. In its launch week, on the Artificial Analysis index then current (v4.1.1), GLM-5.3 scored 60, second of 111 evaluated models, and GLM-5.3-Flash scored 57, fourth, at \$0.09 of cost per task against the flagship's \$1.40, with the best closed model, Claude Opus 5 at maximum reasoning effort, at 63~\cite{glm53-aa} [A]. Those are the numbers the trade press carried and the numbers the shares moved on, and they are also numbers that no longer exist: the re-based v4.3 board of 7~September puts Claude Fable 5.1 and GPT-6 Astra (both at maximum effort) at 53, Claude Opus 5 at 51, Muse Spark 1.3 at 48, GPT-5.6 Sol at 47, \textbf{GLM-5.3 at 45}, \textbf{GLM-5.3-Flash at 42}, Qwen3.8-2.4T-A95B at 40, GPT-5.6 Luna at 38, DeepSeek V4 Pro at 36 and DeepSeek V4 Flash at 35~\cite{aa-v43} [A]. Under the rule of method stated in \S\ref{sec:muse}, the launch-week rank is informative and the launch-week score is not, and the gaps this book prints are the v4.3 ones. Two readings follow, and the second is the one that matters. The open--closed gap at the very top is \emph{eight index points} on v4.3---three on the launch-week board, which is the figure that circulated, and the difference is entirely the instrument: the re-basing replaced Terminal-Bench 2.1 with v4.0, swapped in AutomationBench-AA, dropped GPQA Diamond and raised the private-set weight, and the top of the board re-ordered accordingly. Eight is consistent in direction with the Stanford Index's 2.7\% but is measured on a different instrument with a different scale, so the two numbers corroborate a shape rather than confirming each other, and neither should be read into the other. And the distilled tier is still the interesting one, though the claim has to be made more carefully than the launch-week board invited: in its launch week the 320B/18B model tied GPT-5.6 Sol at high reasoning effort and beat Gemini 3.7 Flash; on v4.3 it sits eleven points below the closed leaders and five below Sol, but above Qwen3.8's 2.4-trillion-parameter flagship, above GPT-5.6 Luna and above both DeepSeek V4 tiers---an artifact less than half the size of its own flagship, three points behind it, at \$0.25 per index task against the flagship's \$2.01~\cite{aa-v43,glm53-flash-agentic} [A]. One caveat belongs in print: the launch-week Flash score circulated in Z.ai's own materials before it appeared on the public board, which is a provenance irregularity rather than a refutation, and this book prints the number with that history attached~\cite{glm53-aa}. Adoption followed at the speed the ecosystem now runs at---346{,}516 Hugging Face downloads in the first month, 62 quantized derivatives, a listing on the Dell Enterprise Hub---and Zhipu's Hong Kong-listed shares rose sharply on the reveal~\cite{glm53-flash,glm53-chips} [V]. What that 320B/18B model does to the hardware requirement is the subject of \S\ref{sec:quant}, and it is the single most consequential measurement in this chapter.

\emph{What the small model is actually good at, stated exactly.} The composite index hides the shape of the result, and the independent v4.3 sub-scores are more useful than the headline. On Terminal-Bench v4.0---the agentic-coding component whose introduction drove the re-basing---GLM-5.3-Flash scores \textbf{33\%}, against 12\% for GPT-5.6 Luna and 12\% for DeepSeek V4 Flash, the two models it competes with on price, and 42\% for its own flagship; on AutomationBench-AA it scores 60\% against Luna's 50\%. On everything that is not agentic coding it loses to the same cheap closed model: 15\% against 24\% on GDP.pdf, 15\% against 21\% on CritPt, 61--71 output tokens per second against Luna's 120, 34.5 seconds to the first answer token, and a cost per index task of \$0.25 against Luna's \$0.18~\cite{glm53-flash-agentic} [A]. The precise finding is therefore not that a distilled Chinese model has caught the closed volume tier; it is that GLM-5.3-Flash is \emph{the cheapest point on the price--agentic-coding frontier}, and not a general replacement for the flagship tier---an 18B-active model that does one economically important thing at a third of the cost of anything that does it better, and that is slower, worse-calibrated and less knowledgeable than the closed model priced beside it at everything else. The vendor's own card claims considerably more, 84.3 on Terminal-Bench 2.1 and 63.4 on DeepSWE v1.1, and neither figure has been replicated by anyone; no SWE-bench Verified, SWE-bench Pro, LiveCodeBench or Aider result for the model exists from any source, vendor or independent, so the coding claim rests on two independent agentic sub-scores and nothing else, and this book says so rather than filling the gap with the model card~\cite{glm53-flash-agentic} [A/C on the vendor figures]. The hardware statement that belongs beside the capability statement is equally narrow: the model does not run on a consumer GPU or on a 32--64\,GB machine; its floor is roughly 100--128\,GB of unified memory at one to three bits, retaining 71--82\% of its unquantized top-1 accuracy on the quantizer's own harness (\S\ref{sec:quant})~\cite{glm53-local} [P]---a single workstation-class machine, not a laptop.

\emph{The domestic-silicon claim is real, narrow, and easy to overstate---so state it narrowly.} Z.ai's own announcement says the model was ``previously previewed as Ox~Alpha, \emph{running} entirely on Chinese AI chips''; founder Jie Tang adds, ``Powered by pure Chinese chips''~\cite{glm53-chips} [P]. Four qualifications are load-bearing and none is a quibble. \emph{First, the claim is about serving, not training}: Z.ai has made no training-hardware statement for GLM-5.3 or for Flash, and implicator.ai's formulation is exact---served is not trained. \emph{Second, no chipmaker is named}: Counterpoint's Ivan Lam speculates Ascend plus other domestic vendors, which is analyst inference and not disclosure, and CNBC reported that it could not verify the claim independently [A]. \emph{Third, it is not a first}: Zhipu stated that GLM-Image was trained end to end on Huawei Atlas hardware in January 2026, with no server count or timing disclosed, and the GLM-5 family was reported in June 2026 to have been trained on roughly 100{,}000 Ascend 910B with MindSpore and no NVIDIA hardware at any stage~\cite{zhipu-ascend-prior}. \emph{Fourth---and this is the qualification that changes the grade of the third---the June claim has no locatable primary.} It circulates through TechTimes, Tom's Hardware and Decrypt without a cited source; the GLM-5 repository carries no hardware statement; and Huawei's own Ascend announcement of February 2026 describes ``0-day'' \emph{inference} support and post-release training guides for GLM-5, which is adaptation after the fact, not a training claim~\cite{glm5-ascend-unsourced}. This book therefore grades the GLM-5 training claim as press-derived [A/C], not as an unaudited company statement, and does not stack anything on it; the only Zhipu domestic-\emph{training} claim with a primary behind it is an image model of undisclosed scale [P, unaudited]. Z.ai's own Chinese-language documentation for Flash is consistent with that reading, because it claims exactly what the English announcement claims and no more: the phrase is \emph{shouci zai daguimo liuliang zhong \ldots\ tigong fuwu}---the first time large-scale traffic has been \emph{served} on a domestic cluster---and it says nothing about training [P]. Against all of it stands the reporting that DeepSeek's R1 training did not complete reliably on Ascend and reverted to NVIDIA~\cite{ascend-delay}---which is precisely why the training/serving distinction carries weight: the one domestic training attempt with independent reporting attached to it is the one that failed. What \emph{is} new, and it is substantial, is production-scale inference: a reported cluster of roughly 100{,}000 domestic accelerators serving 62 trillion tokens over the preview week while carrying a workload that ranked fourth in the world on the launch-week (v4.1.1) board~\cite{glm53-chips} [A]. Circulating reports of ``100 trillion tokens per day'' are irreconcilable with router data and are not used here.

Score it against the test this book actually set. Proposition~P2 is confirmed by a frontier-class model \emph{trained} end to end on domestic silicon, packaging and memory---claim CR-09, the ``Mate~60 moment'' of AI. August 2026 supplies the serving half of that stack at production scale, on a workload that was top-four in the world in its launch week, under a company statement that is unusually specific about what it does \emph{not} say. That is the closest the integration test has come to being passed, and it is not a pass---and the fourth qualification above makes the failure cleaner, not closer: no training claim has been made for these models; the one training claim for a predecessor that has a primary is an image model with no scale disclosed; the frontier-scale training claim that the trade press repeats has no primary at all; no vendor is named; and the memory and packaging inside that cluster are undisclosed. A frontier model trained end to end on domestic silicon, packaging and memory has not been passed as a test, and the evidence that it has is thinner at \datafreeze than the June headlines made it look. CR-09 therefore remains ``not yet observed'' and this book still does not assert it. The revision the evidence does support is narrower and still consequential: the \emph{inference} leg of sovereign compute has moved from roadmap to production, which is the leg that carries the deployed token volume, sets the standards, and pays the bills---and it has done so on the workload the whole world is currently ranking.

\subsection{The barbell: where the tokens go and where the money stays}
\label{sec:barbell}

The router and gateway data of the fortnight before \datafreeze make an argument that the index tables cannot, and it is the spine of this book's reading of the small model. Take the volume side first. On OpenRouter, with data through 8~September, the four largest models by routed tokens were Tencent's Hy4 preview at 18.7 trillion, GPT-5.6 Luna at 14.5 trillion, DeepSeek V4 Flash at 12.4 trillion and \textbf{GLM-5.3-Flash at 12.2 trillion}; the rest of the top ten were DeepSeek V4 Flash 0423, MiniMax M3, Hy3, NVIDIA's Nemotron~3 Ultra, Xiaomi's MiMo-V2.5 and GLM-5.3 itself, at 3.4--5.0 trillion each. Where the active-parameter count is disclosed at all, the top ten runs 49B, 13B, 18B and below: the volume tier is a sparse-mixture tier, and the frontier flagships are not in it. An OpenRouter engineer, quoted in the same week, put Claude Fable~5---of the closed family that tops the v4.3 board---in ``the low billions'' of monthly output tokens on the platform against roughly 12 trillion each for Luna and GLM-5.3-Flash, on a platform whose total volume has risen more than twenty-five-fold in a year and doubled in the last month~\cite{openrouter-sep26} [A]. Router volumes are a small, promotion-distorted slice of global inference and are read here as a direction and not a census; but three orders of magnitude is a direction.

Now the spend side, from a gateway that sees the invoices as well as the tokens. On Vercel's AI Gateway, open-weight models took 36\% of July's tokens on \textbf{8.6\% of spend}, and 62\% of tokens on a single day, 22~August; 81\% of tokens ran on models that did not exist six months earlier; and the gateway puts the self-hosting break-even ``past 50 million tokens per day.'' In the same data the four leading frontier laboratories still take 95\% of spend, and Anthropic alone 61--65\% of it, at roughly 4.4 times the average token price~\cite{vercel-aug26} [A]. That is a barbell: the sparse open tier is winning \emph{volume} and losing \emph{spend}, and the closed frontier is doing the reverse, and the two halves of the market are separating rather than converging. It is also, read against Section~\ref{sec:io}, exactly the rent-relocation pattern this book predicted, arriving in the accounts: the model layer at the volume end has been commoditized to a price at which a gateway's customers move to whichever sparse model shipped last month, and the rent has moved up---into the capability tier, into the harness (Chapter~\ref{ch:harness}), into the hardware, and, as \S\ref{sec:toll-road} shows, into the terms on which the commons is served at scale. Morgan Stanley's formulation of the consequence is the correct one and the correctly limited one: open weights ``may really kill\ldots\ the idea that the foundation model layer can sustain high rents''---the rent, not the compute demand, which the same barbell shows growing on both arms~\cite{slm-belcak} [A].

The small-model argument has a literature, and the literature cuts both ways in the way the data do. Belcak and colleagues at NVIDIA Research argued in 2025 that most agentic sub-tasks are repetitive, narrow and non-conversational, and are therefore better served by small models with a large model held in reserve---a thesis the router data now look like a demonstration of, and one whose author's own laboratory then shipped a 120B agentic flagship, which is not small~\cite{slm-belcak} [A]. Against it stands the ``Hierarchy of Agentic Capabilities'' study of January 2026: on 150 realistic workplace tasks, ``weaker models struggle with fundamental tool use and planning,'' and even the best models fail about 40\% of them~\cite{slm-belcak} [A]. The Flash sub-scores of \S\ref{sec:glm53} are the two findings in one model: a third of Terminal-Bench v4.0 is the best result at its price and a long way from a solved workload, and the composite index still puts it eleven points below the closed leaders. The counter the book has to print is therefore the strongest one available, and it is in the same data as the thesis: spend concentration in the frontier is evidence that the capability tier still commands a rent, that buyers who can measure the difference pay 4.4 times the average price for it, and that the 40\% of workplace tasks nobody's model completes is where that rent is earned. What the barbell does not show is which arm grows faster, because token volume and spend are being measured on different customers---the volume arm on a router whose users self-select for price, the spend arm on a gateway whose largest customers self-select for capability---and no source at \datafreeze reports both on one population. The book's position is correspondingly narrow: the volume tier is Chinese, sparse and open, and no longer commands a rent; the capability tier is American, metered and closed, and still does; and the strategic question of Chapter~\ref{ch:trade} is not whether the first is true but how long the second lasts.

\subsection{AI-designed silicon: why the compute dependence is transitional}
\label{sec:offshore}

The Qwen3.8-Max release supplies what may be the single most consequential datum in this chapter, and it bears directly on the one element of the Chinese stack that Part~II has consistently graded weakest.

Alibaba's own announcement lists, among Qwen3.8-Max's headline capabilities, ``system-level autonomous planning with closed-loop adaptive learning, driving \emph{500+ turns of chip design optimization}''~\cite{qwen-announcement}, and its accompanying material claims the model has ``independently achieved the autonomous execution of the entire silicon design flow''~\cite{berman2026}. The observation that makes this strategic rather than merely impressive was put most directly by an independent Western commentator reviewing the release: \emph{one of the things China is behind on is chip design and chip manufacturing---but a model that can help design those chips means that is no longer the weakness}~\cite{berman2026}.

Set that against the architecture of Part~II and the loop closes with unusual precision. China holds an open instruction set with no revocable licence (Section~\ref{sec:riscv}); domestic logic at 5\,nm-class without EUV; a memory strategy that routes around HBM (Section~\ref{sec:lpddr6}); an open interconnect and software stack; a demonstrated exascale integrated-CPU machine (Section~\ref{sec:lineshine}); and a national EDA push against the one remaining Western duopoly. The binding constraints on that stack have never been ambition or capital---they have been \emph{design throughput and tooling}: the accumulated human expertise and proprietary software that turn an architecture into a manufacturable die. That is exactly the input an autonomous silicon-design flow attacks, and it is the input least protected by export control, because a design capability that a Chinese lab trains cannot be embargoed at a border.

This is the flywheel of Section~\ref{sec:flywheel} operating on the element the book had marked as slowest. AI improving the design of the chips that run the AI is the recursive loop in its most literal form, and it arrives---if the claims survive scrutiny---aimed squarely at the gap between the open-frontier position China already holds and the sovereign-compute position it does not yet. The corollary for this book's evidence base is that the current arrangement, in which Chinese labs reportedly train on NVIDIA capacity leased offshore in Singapore and Malaysia while serving domestically~\cite{qwen-offshore}, should be read as \emph{transitional rather than structural}. It is what a fast-moving ecosystem does while its own silicon matures; it is not a stable dependence, and the co-design mechanism that ends it is now visible in a shipped product rather than a roadmap. Claim CR-09---a frontier-class model trained end-to-end on a domestic stack---remains ``not yet observed,'' and this book still does not assert it; \S\ref{sec:glm53} shows how far the \emph{serving} half of that stack has come and why the training half is still not in evidence. But the pathway to it has shortened, and the forecast register is revised accordingly.

The most striking corroboration comes from a commentator who reaches the book's conclusion while explicitly not wanting to. Reviewing the release for a Western audience, Berman traces the adoption logic---US enterprises ``can get 95\% of the capability and pay a fraction of the price,'' and for the vast majority of use cases ``the absolute frontier just isn't necessary and in fact might be overkill,'' so they will choose the open models, which are Chinese---and then follows it one step further to the conclusion Part~II of this book is built on: \emph{``eventually what will likely happen is the model will be co-designed with the chip to be the most efficient and the most optimized for that specific chip\ldots{} with extreme co-design between model and hardware it might actually put us, as we're becoming dependent on Chinese open source, highly dependent on Chinese chips as well''}~\cite{berman2026}. That is the convergence thesis of this book, stated independently, by an observer who calls the resulting geopolitical exposure ``not my opinion---it is just a fact,'' and who is manifestly unhappy about it. When the mechanism this book has spent Part~II describing is arrived at independently by someone who would prefer it were false, the argument has stopped being idiosyncratic.

Two disciplines keep this section honest. \emph{The claims are the vendor's}: 500 turns of chip-design optimization and an autonomous silicon flow are Alibaba's descriptions of Alibaba's model, and the appropriate verification event---an independently reproduced tape-out-grade design, or third-party evaluation of the flow---has not occurred. The GitHub project trace published alongside the autonomous-coding demonstration is a genuine step toward auditability and more than most Western frontier claims carry; the chip-design claim has no equivalent artifact yet. \emph{And the counter-argument is real}: the same review notes that the leading US closed models are rumoured at roughly seven trillion parameters, against 2.4--2.8T for the Chinese frontier---roughly twice the scale, which requires the best accelerators in the world and a great many of them~\cite{berman2026}. The Chinese achievement is therefore properly read as \emph{capability parity at a third of the parameters and a fraction of the price}, which is an efficiency victory rather than a scale victory---and efficiency victories are exactly what Part~I's precedents are made of, but they are not the same thing as removing the compute constraint. Section~\ref{sec:rsi-branch} takes the strongest form of that counter-argument seriously.

Kimi K3 deserves one more paragraph, because it is the cleanest single demonstration of the five-frontier framework of \S\ref{sec:scoreboard}. Its own technical report confirms 2.8T total parameters, 104B activated, a one-million-token context, native vision---and explicitly acknowledges that the model still trails the strongest proprietary systems overall~\cite{kimiK3-report}. Here is the open frontier and the absolute frontier in a single primary source: the largest, most capable open-weight artifact ever released, self-described as behind the closed leaders. Both of this book's core observations---that China leads the open frontier decisively, and that the absolute frontier has not dissolved---are stated by the releasing laboratory itself.

\textbf{Alibaba (Qwen)} is the adoption champion: the most-downloaded open family on Earth, crossing one billion cumulative Hugging Face downloads in July 2026 with over 113,000 derivative models---more than Google and Meta combined, and roughly 40\% of all new LLM derivatives on the Hub~\cite{hf-spring-2026,qwen-billion}. \textbf{Moonshot (Kimi)} is the scale champion: K3's 2.8 trillion parameters (896 experts, $\sim$1.8\% activation, hybrid linear attention, native MXFP4, 1M context) make it the largest open-weight model ever released, ranking in the global top five overall---the only open model in that tier~\cite{kimi-k3}. \textbf{DeepSeek} is the efficiency champion, each generation shipping an attention or sparsity innovation the rest of the field adopts within months. \textbf{Zhipu (GLM)}---which listed in Hong Kong in January 2026 as ``the world's first publicly traded foundation-model company''---holds the open-weight quality crown on independent indices, at 45 on the re-based v4.3 Artificial Analysis index against 53 for the best closed systems (60 against 63 on the launch-week v4.1.1 board), and is reported in the trade press---with no primary that this book could locate---to have trained GLM-5 on Ascend, a claim graded [A/C] and separated carefully in \S\ref{sec:glm53} from the serving claim Z.ai actually makes for GLM-5.3-Flash~\cite{glm5,glm53-aa,aa-v43,glm5-ascend-unsourced,ipos}. Around them: \textbf{MiniMax} (M-series, agentic coding), \textbf{StepFun} (throughput-optimized), \textbf{Xiaomi} (MiMo, which took the largest single-provider share on OpenRouter within months of entry), \textbf{ByteDance} (Doubao at 345M MAU and $>$120 trillion tokens/day; China's top-rated Arena model), and \textbf{Tencent Hunyuan}~\cite{digitalapplied,datagravity}. Zhipu and MiniMax IPO'd in Hong Kong in January 2026; Moonshot raised at a \$20B valuation~\cite{ipos}, and by September had filed confidentially for a Hong Kong listing of its own at a reported pre-money valuation near \$50B~\cite{moonshot-ipo} [A]. DeepSeek, which had taken no outside capital at all, took its first in June 2026: roughly \$7.4B at a valuation above \$50B, with the state's National AI Industry Investment Fund---the ``Big Fund''---as the sole holder of direct equity and voting rights and Tencent and CATL in a non-voting partnership behind it~\cite{deepseek-round} [A]; a later report of a \$74B valuation conflicts with that account and is not used here. The LLM price war is the ``531'' cull in progress: token prices at 1/7th to 1/150th of US equivalents, margins near zero, survivors selected by capability-per-yuan~\cite{token-econ}.

\section{The toll road: the licence becomes a claim on the clouds}
\label{sec:toll-road}

Section~\ref{sec:maxtier} ended with a licence clause and a prediction: that rent, driven out of the weights by the ratchet, would relocate to the terms of large-scale service, and that the releasing laboratory would retain a claim there. In the last fortnight of August the claim acquired a counterparty, and the counterparty is the American hyperscaler.

\emph{The fact pattern, at the strength the sourcing supports.} On 26~August Reuters reported, on three anonymous sources, that Moonshot AI is in \emph{early-stage} talks with Microsoft, Amazon and Google over agreements under which the three clouds would host Kimi K3 and Moonshot would take a share of the revenue---the laboratory seeking \textbf{up to 30\% of K3-related service revenue}, on terms said to match those offered to its large direct customers. The revenue split, data access and the auditing of token counts were all unresolved; the three clouds declined comment and Moonshot did not respond; nothing was signed, and Reuters' own formulation was that there is ``no certainty'' agreements will result. Smaller platforms had already signed: Chinasoft International disclosed a ``token revenue-sharing and joint innovation agreement'' with Moonshot on 20~July, percentages undisclosed~\cite{moonshot-cloud-talks,kimi-k3-license-maas} [A]. Every account says \emph{revenue}, not profit, and none attributes to the talks any motive beyond monetization; this book does not add one.

\emph{What forces the negotiation is the licence, not diplomacy.} The reason a hyperscaler must talk to Moonshot at all is a clause that did not exist for K2. The K3 licence provides that a licensee operating a model-as-a-service business whose aggregate revenue exceeds \$20 million over any consecutive twelve months ``must enter into a separate agreement with Moonshot AI before using the Software or its derivative works for any commercial purpose,'' with a display requirement above 100 million monthly users or \$20 million of monthly revenue~\cite{kimi-k3-license-maas} [P]. That is consistent with K3, unlike K2.5, not being on Amazon Bedrock, and why it is reachable in Microsoft Foundry only through Fireworks AI as an intermediary partner model, at \$3.30/\$16.50 per million tokens~\cite{us-clouds-chinese-weights,kimi-k3-license-maas} [P]. Alibaba is copying the instrument: Reuters reported on 7~August that large commercial users of Qwen3.8-Max will be required to share revenue, the rate under negotiation and the design ``modelled on Moonshot's,'' with the \$50 million threshold of the \texttt{qwen3.8-max} licence as the lever~\cite{qwen-revshare} [A]. The clause is the mechanism; the hyperscaler is the counterparty; and the rent Section~\ref{sec:io} said would leave the weights has arrived, in draft, in a contract.

\emph{Against what already exists.} The talks are startling only against the baseline, so the baseline has to be stated from the clouds' own documentation. Amazon Bedrock has served DeepSeek-R1 since January 2025 and, since February 2026, DeepSeek V3.2, MiniMax M2.1, GLM~4.7 and 4.7 Flash, Kimi K2.5 and Qwen3 Coder Next as fully managed models; Google Vertex serves Kimi K2 Thinking, MiniMax M2, DeepSeek-V3.2, GLM~4.7 and GLM~5; Microsoft Foundry lists DeepSeek in its ``sold by Azure'' category. AWS's third-party model terms state that DeepSeek and Qwen serverless models ``are sold by AWS'' and that customer data is not shared with the model providers. No cloud discloses any payment to any Chinese laboratory for serving permissively licensed weights, and Bedrock lists DeepSeek V3.2 at roughly two to five times the laboratory's own price~\cite{us-clouds-chinese-weights} [P]. The American hyperscalers, in other words, are already the largest Western resellers of the Chinese commons, at a markup, for nothing. What the Moonshot talks would be---if they close---is the first time the maker of the commons takes a cut of the commons' service revenue on a Western cloud: a toll booth built on somebody else's road, and the sharpest instance yet of the result Section~\ref{sec:io} predicted, because the weights remain free and the rent is collected one layer up, from the party with the customers.

\emph{The ``hedge'' reading, stated and rejected on the evidence.} It is tempting to read the talks as insurance---a Chinese laboratory buying protection against a Western shutdown by making American clouds its partners, or the clouds buying protection against a Chinese one by contracting for supply. No source attributes either motive, and the direction of the legal exposure runs the other way. Webster and Xu's analysis of the K3 licence is the only serious treatment located, and its title is its argument: \emph{with revenue comes regulability}. Downloading open weights left it genuinely ambiguous whether any transaction with a Chinese laboratory had occurred; ``if a US company needs a contract with Moonshot to provide the inference tokens,'' the picture changes, under a supply-chain rule ``it's harder to argue there's no transacting going on,'' and if Moonshot were sanctioned ``the contract to provide inference is clearly an instance of doing business''~\cite{webster-regulability} [A]. The instruments that would give that exposure teeth are partly enacted and partly pending. Section~1532 of the FY2026 National Defense Authorization Act, in force since December 2025, requires the Department of Defense to exclude ``covered AI''---defined by reference to DeepSeek and High-Flyer by name, to the Consolidated Screening and 1260H lists, and to covered-nation domicile---and does not distinguish self-hosted weights from cloud-served tokens; the No Adversarial AI Act remained in committee at \datafreeze~\cite{ndaa-1532} [P]. Treasury Secretary Bessent was reported in July to have threatened Entity-List designation of Chinese laboratories [A]; Axios reported the same month a draft executive order that would make US hosts of Chinese models guarantee their security and accept liability, not issued at the currency date~\cite{webster-regulability} [A]. And on \datafreeze itself the NSA, FBI and CISA issued a joint advisory naming DeepSeek, Moonshot, Alibaba, MiniMax, StepFun and Z.ai as running distillation campaigns against US frontier models, with the traffic routed through cloud providers and third-party aggregators---an advisory about distillation rather than hosting, but one that names the prospective toll-road partner and the clouds in the same paragraph~\cite{cisa-distillation} [P]. A revenue contract with a US hyperscaler is therefore not a shield against any of that; it is the document that would convert an ambiguous download into an unambiguous transaction. On the Chinese side the picture is the mirror image and equally unresolved: the Ministry of Commerce spent June and July consulting Alibaba, ByteDance and Z.ai on options that include barring public release of the most advanced models, restricting overseas access and criminalizing model leakage, with ``nothing decided'' and no timeline, and Bloomberg reported in May that passport-surrender rules first applied at DeepSeek had been extended to Moonshot, StepFun and ByteDance~\cite{mofcom-model-curbs} [A]. If there is a shutdown the toll road hedges against, then, it is the Chinese one---a contracted revenue stream from Redmond, Seattle and Mountain View being a harder thing for Beijing to close than a download link. That is this book's inference and not anyone's reporting, and it is graded accordingly [S].

\emph{What it does to the thesis.} The commons is acquiring a cash register, and the two facts that matter are that it remains a commons and that the register sits exactly where the theory said it would. Below the threshold nothing changes: every researcher, every enterprise deploying internally, every derivative-builder and every service business under \$20 million of revenue uses K3 on the terms it was released under, and the download counts of Table~\ref{tab:metric-matrix} continue to accrue. Above it, the same artifact is a toll road, and the toll is levied not on the user but on the reseller with the largest customer base on earth. That is ``openness relocates rent'' in its most literal form yet, and it is the result of Section~\ref{sec:io} and \S\ref{sec:maxtier} with a named counterparty and a percentage attached. The counter is the size of the register. Moonshot's reported annualized revenue was \$300 million in June, with no geographic split disclosed, against a frontier laboratory's run rate two orders of magnitude larger~\cite{moonshot-ipo} [A]; 30\% of a hyperscaler's K3 line is a claim on a market that NVIDIA's Poolside purchase and its Nemotron line are built to commoditize from the other side (Chapter~\ref{ch:mirror}); and a negotiation that Reuters describes as early-stage, unsigned and silent on its three hardest terms is evidence of intent, not of income. The toll road is real as a mechanism and unproven as a business, and the ratio between those two statements is the ratio this chapter has found at every layer of the model element.

\section{Three scoreboards, not one}
\label{sec:scoreboard}

Before the numbers, a distinction that resolves most of the argument about whether China has ``caught up.'' There are \emph{five} different frontiers, and conflating them produces both triumphalism and complacency:

\begin{description}\itemsep2pt
  \item[1. The absolute frontier] --- the highest demonstrated capability under controlled evaluation, irrespective of access or price. The United States leads this, and the lead is real.
  \item[2. The open frontier] --- the highest capability available as downloadable weights under usable licence terms. China leads this, and not narrowly.
  \item[3. The deployable frontier] --- the highest capability an institution can \emph{operate reliably at acceptable cost}, after quantization, serving constraints, security review, and licence compatibility. China increasingly leads this.
  \item[4. The economic frontier] --- the system with the lowest cost per \emph{successful task} for a given workload (\S\ref{sec:cost-solved}). This is frequently neither principal's flagship: for high-volume routine work it is often an older, smaller model that nobody writes about.
  \item[5. The scientific frontier] --- the system that produces the greatest validated scientific value per unit of national resource. This is the frontier Chapter~\ref{ch:prize} argues matters most, and the one nobody currently measures.
\end{description}

\begin{figure}[htbp]
\centering
\begin{tikzpicture}[
  every node/.style={font=\scriptsize},
  fr/.style={draw=inkMut, line width=0.6pt, rounded corners=2pt, align=left,
             inner sep=4pt, text width=42mm, minimum height=10mm},
  pcell/.style={align=center, font=\scriptsize, text width=22mm, inner sep=2.5pt,
              rounded corners=2pt, minimum height=6mm},
]
\node[font=\scriptsize\bfseries\color{inkPri}] at (4.35,0.75) {leader today};
\node[font=\scriptsize\bfseries\color{inkPri}] at (7.05,0.75) {decides what};

\node[fr, fill=sOne!12, draw=sOne] (f1) at (0,0) {\textbf{1. Absolute}\\ best under controlled evaluation};
\node[pcell, fill=sOne!22] at (4.35,0) {United States};
\node[pcell, fill=surf] at (7.05,0) {prestige, hardest tasks};

\node[fr, fill=sThree!12, draw=sThree, below=3.5mm of f1] (f2) {\textbf{2. Open}\\ best downloadable under usable licence};
\node[pcell, fill=sThree!22] at (4.35,-1.4) {China};
\node[pcell, fill=surf] at (7.05,-1.4) {developer default};

\node[fr, fill=sThree!12, draw=sThree, below=3.5mm of f2] (f3) {\textbf{3. Deployable}\\ best operable reliably at acceptable cost};
\node[pcell, fill=sThree!22] at (4.35,-2.8) {China (rising)};
\node[pcell, fill=surf] at (7.05,-2.8) {industrial outcomes};

\node[fr, fill=sTwo!12, draw=sTwo, below=3.5mm of f3] (f4) {\textbf{4. Economic}\\ lowest cost per \emph{successful task}};
\node[pcell, fill=surf] at (4.35,-4.2) {often neither: an older, smaller model};
\node[pcell, fill=surf] at (7.05,-4.2) {the majority of tokens};

\node[fr, fill=sFour!12, draw=sFour, below=3.5mm of f4] (f5) {\textbf{5. Scientific}\\ validated scientific value per unit of national resource};
\node[pcell, fill=surf] at (4.35,-5.7) {unmeasured};
\node[pcell, fill=sFour!22] at (7.05,-5.7) {long-run national power};

\node[font=\scriptsize\color{inkSec}, align=left, text width=146mm, anchor=north west] at (-2.3,-6.6)
  {Apparent contradictions dissolve here. China can trail frontier~1 while leading 2 and 3. A US closed system can
   dominate frontier~5 for a given workload despite far higher token prices, because scientific value per resource
   depends on capability at the hard end rather than price at the volume end. And frontier~4 is frequently defined
   by a model nobody writes about.};
\end{tikzpicture}
\caption[The five frontiers]{Five frontiers, won separately. Only the first two are systematically measured; the
third and fourth are measurable but rarely reported (\S\ref{sec:cost-solved}); the fifth---the one
Chapter~\ref{ch:prize} argues matters most for national outcomes---has no accepted metric at all, which is both a
gap in the literature and an invitation to whoever builds the first credible one. Positions are this book's
assessment as of the version stamp [S].}
\label{fig:frontiers}
\end{figure}


Many apparent contradictions dissolve under this framework (Figure~\ref{fig:frontiers}). China may trail the absolute frontier while leading the open and deployable ones. A US closed system may dominate the \emph{scientific} frontier for a particular workload despite far higher token prices, because scientific value per resource depends on capability at the hard end rather than on price at the volume end. And an eighteen-month-old open model may define the economic frontier for the majority of tokens actually generated in the world. Frontier~3 is where the two halves of this chapter meet: a model that ranked fourth in the world in its launch week (v4.1.1) and sits eleven index points below the closed leaders on the re-based v4.3 board, and that fits, quantized, in the memory of one workstation (\S\ref{sec:glm53}, \S\ref{sec:quant}) is a capability an institution can operate without a vendor relationship, a network egress, or a procurement fight, which is a different kind of lead from the one the index tables measure. That China may lead frontiers 2, 3, and 4 without leading 1 \emph{strengthens} this book's thesis rather than weakening it, because industrial victory in every precedent of Part~I depended on what a buyer could actually install, at what price, under what terms---not on the laboratory maximum. The rest of this section reports frontiers 1 and 2, which are measured; Chapters~\ref{ch:trust} and~\ref{ch:sizing} supply what is known about 3 and 4; frontier 5 is, at present, an open measurement problem and an invitation.

Three independent measurement families bear on the adoption proposition (P1), and they must be kept apart, because they measure different things and currently disagree in one important place. \emph{Family one---capability}: index and arena measurements of model quality. \emph{Family two---distribution}: downloads, derivatives, and routed tokens. \emph{Family three---production and revenue}: enterprise deployments and paid spend. Families one and two tell one story (Table~\ref{tab:model-scoreboard}); family three lags it, and honesty requires displaying the lag.

On capability: the Stanford AI Index puts the top US--China gap at 2.7\% as of March 2026, down from double digits in 2023---while also finding that the closed-over-open gap \emph{reopened} to 3.3\% (from 0.5\% in August 2024) and that six of the top ten arena models are closed~\cite{stanford-aiindex-2026}. Both facts are true at once: China reached the top competitive tier; the absolute closed frontier did not dissolve. A second instrument agrees on the shape and, at \datafreeze, disagrees with itself on the size: on the Artificial Analysis index the best open-weight Chinese model, GLM-5.3, scores 45 on the re-based v4.3 board against 53 for the best closed models, Claude Fable 5.1 and GPT-6 Astra---\emph{eight} index points---where the launch-week v4.1.1 board of a fortnight earlier had it at 60 against 63, three points, and the 320B/18B GLM-5.3-Flash at 42 on v4.3 sits eleven behind the leaders, having been fourth in the world and level with GPT-5.6 Sol on the board it launched on~\cite{glm53-aa,aa-v43} [A]. The Stanford figure and the Artificial Analysis figure are different instruments on different scales, and a 2.7\% gap on one is not three points or eight points on the other; two instruments agreeing that the gap is small and closed-favouring is corroboration of a direction, not confirmation of a magnitude, and neither instrument measures the hard agentic tail where the closed lead is widest---which is exactly the tail the September re-basing weighted more heavily, and the re-ordering it produced is itself evidence of where that lead lives. Benchmark figures additionally carry protocol variance that vendor tables hide---Kimi's own documentation shows 65.8\% on SWE-bench Verified single-attempt becoming 71.6\% with parallel test-time compute~\cite{kimi-k2-thinking}, and reasoning-model response lengths are inflating $\sim$5$\times$ per year [V]~\cite{epoch-output}, so a model ten times cheaper per token that spends ten times the tokens per solved task has gained nothing per task. On distribution and production, Table~\ref{tab:metric-matrix} states what each metric can and cannot support.

\begin{table}[htbp]
  \centering
  \small
  \caption{The adoption-metric matrix: what each measure proves.}
  \label{tab:metric-matrix}
  \begin{tabular}{p{38mm}p{50mm}p{54mm}}
    \toprule
    Metric & Measures & Principal limitation \\
    \midrule
    HF downloads (41\% Chinese) & Distribution, developer interest & Automation/CI inflation; no proof of production use~\cite{hf-spring-2026} \\
    Derivatives ($>$113k Qwen) & Ecosystem reuse & Many experimental or dormant \\
    OpenRouter tokens ($\sim$61\% mid-2026) & Usage on one neutral router & Platform-specific mix; not a census~\cite{openrouter-2026} \\
    API revenue & Monetized hosted demand & Misses self-hosting entirely \\
    Enterprise deployment (DeepSeek $>$26k orgs [P]) & Operational adoption & Vendor-survey selection~\cite{sos-ai} \\
    Enterprise \emph{spend} share & Follows the money & Fell 19\%$\rightarrow$11\% in 2025---the countertrend~\cite{menlo} \\
    \bottomrule
  \end{tabular}
\end{table}

One adoption channel none of those six metrics registers is the silicon industry's own. At Hot Chips 2026, nine separate decks reported inference economics on Chinese open-weight models: d-Matrix on GLM~5.2 and Kimi~K3 serving rates; OXMIQ building its cost model on Kimi-K2; NVIDIA measuring agentic throughput on DeepSeek-V4-Pro in the Rubin deck and using DeepSeek~V3 training-step time as the headline multi-tenancy result in the Spectrum-X deck; Intel plotting every published frontier model from Llama~2 70B to Kimi~K2 to establish its capacity-versus-bandwidth thesis; SambaNova quoting MiniMax M2.7; OpenAI reporting DeepSeek~R1 on GB300~\cite{hc2026-china-bench} [V]. No Chinese chip was named anywhere in the conference and no Chinese vendor presented; what the Chinese laboratories supplied was the \emph{workload}. American silicon vendors benchmark and publish against Chinese weights because those are the frontier-scale weights they can obtain, run, and quote in public---which makes Chinese model artifacts a de facto reference load in the design and marketing of Western accelerators, a form of standard-setting that no download counter captures and that runs in the opposite direction to the procurement bans of Chapter~\ref{ch:trust}.

\begin{table}[htbp]
  \centering
  \caption{The model element, mid-2026 snapshot.}
  \label{tab:model-scoreboard}
  \begin{tabular}{p{72mm}p{72mm}}
    \toprule
    Capability & Adoption \\
    \midrule
    US--China top-model gap: 2.7\% (Stanford AI Index 2026), from 17.5--31.6\% in 2023~\cite{stanford-aiindex-2026} &
    Chinese models: 41\% of all Hugging Face downloads (trailing year); China passed the US in cumulative downloads~\cite{hf-spring-2026} \\
    \addlinespace
    Release lag behind US frontier: from 6--8 months to ``a matter of weeks''~\cite{aei-2026} &
    OpenRouter routed tokens: Chinese models $\sim$2\% (early 2025) $\rightarrow$ $\sim$61\% (mid-2026); Meta/Llama below 1\%~\cite{openrouter-2026,datagravity} \\
    \addlinespace
    Top open-weight models: GLM-5.3 and GLM-5.3-Flash, Kimi K3, DeepSeek V4, MiniMax M3---all Chinese; the best US open entry is still an NVIDIA model, sixth on the router at 5.62T weekly tokens (Chapter~\ref{ch:mirror})~\cite{openrouter-2026,openrouter-aug26} &
    a16z: of startups pitching with open-source stacks, $\sim$80\% use a Chinese model; DeepSeek counts $>$26,000 enterprise deployments~\cite{casado,sos-ai} \\
    \bottomrule
  \end{tabular}
\end{table}

\begin{figure}[htbp]
\centering
\begin{tikzpicture}
\begin{axis}[bookaxis, height=5.2cm,
  xlabel={}, ylabel={Share of routed tokens (\%)},
  title={Where the tokens went: OpenRouter author share},
  xmin=0, xmax=6, ymin=0, ymax=100,
  xtick={0,1,2,3,4,5,6},
  xticklabels={{Jan\\'25},{Apr\\'25},{Jul\\'25},{Oct\\'25},{Jan\\'26},{Apr\\'26},{Jul\\'26}},
  xticklabel style={align=center, font=\scriptsize\color{inkSec}},
  legend style={at={(0.98,0.60)}, anchor=east},
]
\addplot[color=sTwo, mark=square*, mark size=1.5pt] coordinates {
  (0,90) (1,84) (2,74) (3,62) (4,52) (5,44) (6,35)};
\addlegendentry{US-origin models}
\addplot[color=sOne, mark=*, mark size=1.5pt] coordinates {
  (0,2) (1,7) (2,16) (3,28) (4,40) (5,52) (6,61)};
\addlegendentry{Chinese-origin models}
\node[font=\scriptsize\color{inkSec}, anchor=south east] at (axis cs:6,63) {61\%};
\node[font=\scriptsize\color{inkSec}, anchor=north east] at (axis cs:6,33) {35\%};
\draw[dashed, color=inkMut, line width=0.5pt] (axis cs:5.0,0) -- (axis cs:5.0,100);
\node[font=\scriptsize\color{inkSec}, anchor=south, align=center] at (axis cs:5.0,78) {crossover\\Jun 2026};
\end{axis}
\end{tikzpicture}

\vspace{1mm}

\begin{tikzpicture}
\begin{axis}[bookaxis, height=4.4cm,
  xlabel={}, ylabel={Measured gap (\%)},
  title={Two gaps moving in opposite directions},
  xmin=-0.3, xmax=3.3, ymin=0, ymax=35,
  xtick={0,1,2,3},
  xticklabels={2023,2024,2025,{Mar 2026}},
  xticklabel style={font=\scriptsize\color{inkSec}},
  legend style={at={(0.98,0.98)}, anchor=north east},
]
\addplot[color=sOne, mark=*, mark size=1.6pt] coordinates {(0,31.6) (1,10.0) (2,4.5) (3,2.7)};
\addlegendentry{top US model over top Chinese model}
\addplot[color=sFour, mark=triangle*, mark size=2.0pt] coordinates {(1,0.5) (2,1.8) (3,3.3)};
\addlegendentry{top closed model over top open model}
\node[font=\scriptsize\color{inkSec}, anchor=south west] at (axis cs:2.55,2.9) {2.7\%};
\node[font=\scriptsize\color{inkSec}, anchor=north west] at (axis cs:2.55,3.1) {3.3\%};
\end{axis}
\end{tikzpicture}
\caption[Adoption share and capability gaps]{The model element in two pictures. \emph{Above}: share of tokens
routed through a neutral aggregator by author origin; Chinese-origin models passed US-origin models in June
2026, having been at $\sim$2\% eighteen months earlier [V, platform-specific]~\cite{openrouter-2026,datagravity}.
\emph{Below}: the two gaps the Stanford AI Index measures. The national gap closed to 2.7\%; the closed-over-open
gap \emph{re-opened} to 3.3\% [V]~\cite{stanford-aiindex-2026}. Both facts are true, and the book's argument
requires holding them together: China reached the top competitive tier, and the absolute closed frontier did not
dissolve. Intermediate points are interpolated between reported measurements.}
\label{fig:model}
\end{figure}


Two further facts frame the table. On hard, contamination-resistant agentic benchmarks the closed-open gap runs larger than headline indices suggest~\cite{deepseek-v4}; and the investment asymmetry is staggering---US private AI investment ran \$286B in 2025 against China's \$12B (a figure that understates Chinese state capital, per the Index's own caveat)~\cite{stanford-aiindex-2026}: a 23:1 private-spending ratio buying a 2.7\% measured gap. Recall the solar lesson: Sharp and Q-Cells also outspent Suntech per watt of R\&D. The question is not who spends most for the peak; it is who owns the volume, the derivatives, the developers, and the price floor. Meta's Llama---the West's open standard-bearer---collapsed below 1\% of routed volume and retreated toward closed development; OpenAI's answer, gpt-oss (August 2025), was outclassed within weeks~\cite{gpt-oss}. The family-three lag---spend share falling even as usage share soared---is the strongest current evidence \emph{against} P3, and Chapter~\ref{ch:trade} treats it as such; the next section explains the mechanism by which the curves are nonetheless converging.

\section{Quantization: frontier models on commodity hardware}
\label{sec:quant}

The most under-analyzed force multiplier of the open-weight strategy is quantization. Open weights alone democratize \emph{access}; quantization democratizes \emph{operation}. The claim must be stated with more care than it usually is. For selected models, quantizers, and benchmark suites, community tooling has reduced storage by 60--90\% while preserving much of benchmark performance; degradation is highly task-, layer-, context-, and calibration-dependent, and becomes material at extreme 1--2-bit precision. What is robust is narrower and still consequential: 4-bit has become a \emph{release} format rather than a lossy afterthought, and 2-bit-class dynamic quants retain enough capability to be operationally useful on hardware an order of magnitude cheaper. Any serious evaluation reports a \emph{quality vector}---mathematics, coding, multilingual reasoning, long-context retrieval, tool use, calibration, refusal behaviour, safety---alongside memory savings, throughput, and energy, rather than a single retention number.

\subsection{The dynamic-quantization breakthrough}

The emblematic actor is Unsloth, an eight-person company whose ``dynamic'' GGUF quantizations allocate precision non-uniformly---MoE expert layers driven to 1--2 bits, attention and dense layers held at 4--6 bits, calibrated on curated token sets against KL divergence. Days after R1's release, Unsloth's 1.58-bit dynamic quant cut the 671B model from 720\,GB to 131\,GB (an 80\% reduction) while a naive 1.58-bit quant produced gibberish; the dynamic version retained $\sim$70\% of a demanding code-generation benchmark and the 2.2-bit version 92\%~\cite{unsloth-r1}. A precision discipline applies to all such figures: retention is task-dependent, and at 1--2 bits it degrades unevenly across mathematics and code, factual recall, long-context retrieval, agentic tool reliability, multilinguality, and safety behavior---the numbers quoted here are specific quants on specific benchmarks [P], not a general law that ``90\% compression costs single digits.'' The robust claims are narrower and sufficient: 4-bit is now effectively a release format (below), and 2-bit-class dynamic quants preserve enough capability to be operationally useful on hardware an order of magnitude cheaper. The method scaled with the frontier: Kimi K2 (1T) to 245\,GB at 1.8 bits; K2 Thinking to 245\,GB at 1 bit with ``$\sim$85\% accuracy retention''; GLM-5 (744B) to 176--241\,GB; Kimi K3 (2.8T, 1.56\,TB native) to 594\,GB at $\sim$79\% retention and 861\,GB at $\sim$90\%~\cite{unsloth-models}. Unsloth alone counts $\sim$10 million monthly model downloads~\cite{unsloth-yc}; the broader quantization stack---llama.cpp/GGUF, AWQ (MLSys 2024 best paper), GPTQ, EXL3, bitsandbytes/QLoRA, vLLM's quantized serving matrix---is now standard infrastructure~\cite{awq,vllm-quant}.

Crucially, the labs closed the loop by shipping \emph{natively} quantized frontier weights: DeepSeek in FP8 since V3, Kimi K2 Thinking in INT4 via quantization-aware training (2$\times$ speedup at near-FP16 quality), K3 in MXFP4, OpenAI's gpt-oss in MXFP4~\cite{kimi-k2-thinking,kimi-k3,gpt-oss}. Four-bit inference is no longer a lossy afterthought; it is the release format, co-designed with training. This matters doubly for Chapter~\ref{ch:compute}: native low-precision weights quarter the memory-capacity and bandwidth demand of serving---which is exactly the resource the LPDDR-based architecture supplies in abundance.

\subsection{The commodity-hardware frontier}

The consequence is that frontier-class inference now runs on hardware measured in single-digit thousands of dollars. The documented data points span three classes. \emph{CPU servers}: a \$2,000 used-EPYC build (512\,GB DDR4) runs the full 671B R1 at 3.5--4 tok/s; KTransformers (Tsinghua, SOSP 2025) reached 286 tok/s prefill and 12+ tok/s decode for R1/V3 on one RTX~4090 plus a dual-Xeon board with 1\,TB of DRAM---a few-thousand-dollar machine serving the model that crashed NVIDIA's market cap~\cite{ktransformers,epyc-rig}. \emph{Unified-memory desktops}: a single Mac Studio M3 Ultra (512\,GB unified at 819\,GB/s, $\sim$\$10k) runs R1 4-bit at 17--18 tok/s under 200\,W; two of them ran Kimi K2.5 (1T, native INT4) at 22 tok/s; four in an RDMA-over-Thunderbolt cluster ($\sim$\$40k, 1.5\,TB unified) served K2 Thinking at 25 tok/s under 450\,W~\cite{mac-frontier}. \emph{Dedicated local-AI boxes}: NVIDIA's DGX Spark (128\,GB unified, \$4,700) and AMD's Strix Halo (128\,GB, \$2,300) both run 120B-class models at 34--39 tok/s, and 30+ laptop models ship with local-AI accelerators in fall 2026~\cite{spark-halo}. llama.cpp's MoE CPU-offload flag (August 2025) made the pattern general: one consumer GPU for attention, cheap system DRAM for the experts~\cite{llamacpp-moe}.

Note what this hardware list is, architecturally: \emph{capacity-rich, bandwidth-modest memory attached to modest compute}---the LPDDR6X design point of Section~\ref{sec:lpddr6}, discovered independently and bottom-up by the enthusiast market. The sparse-MoE model family and the capacity-first hardware family are co-evolving at every scale from the \$2,000 desktop to the national cluster.

The sharpest instance yet is GLM-5.3-Flash (\S\ref{sec:glm53}), and it is worth quoting at measurement precision because the temptation to overstate it is considerable. Unsloth's published quantizations of the 320B/18B model run from UD-IQ1\_S at 93.1\,GB, retaining 70.9\% top-1 accuracy against the unquantized reference, through UD-Q2\_K\_XL at 108.7\,GB and 78.3\%, to \textbf{UD-IQ3\_XXS at 120.4\,GB and 81.6\%} and UD-Q4\_K\_XL at 199.7\,GB and 92.2\%; the BF16 weights are 641.6\,GB. Minimum combined RAM plus VRAM is 100\,GB at one bit, 128--150\,GB at three bits and 162--210\,GB at four; Unsloth states that the three-bit quantization ``works on 128GB devices like a Mac or NVIDIA DGX Spark,'' and reports measured decode on four DGX Spark units at native FP8 of 28--34 tok/s on short prompts and $\sim$43 tok/s at 50K-token prompts~\cite{glm53-local} [P]. The comparison that gives the numbers their force is the flagship: full GLM-5.3 at FP8 requires eight H200- or H20-class accelerators. Three disciplines keep the claim honest. This is \emph{workstation and small-server} territory---a single 128\,GB unified-memory box of the class priced above at \$2,300--\$4,700 for the three-bit quantization, and the four-Spark configuration that produced the measured decode rates is a $\sim$\$19{,}000 machine---and emphatically not a laptop; the accuracy figures are Unsloth's own top-1 retention measurements on their own harness, not an independent evaluation across the quality vector this section insists on; and there is no verified throughput measurement on a Mac Studio for this model, so the memory footprint is the claim and the four-Spark decode rate is the only throughput number printed here. With those caveats stated, the strategic content survives all of them: a model that ranked fourth in the world in its launch week (v4.1.1) and stands eleven index points from the closed leaders on the re-based v4.3 board, above the closed volume tier it is priced against, now fits in the memory of a machine a single engineering team can requisition without a procurement cycle---which is the \emph{deployable} frontier of \S\ref{sec:scoreboard} arriving, in a specific configuration, at a specific measured size.

\subsection{The metric that matters: cost per solved task}
\label{sec:cost-solved}

Price per million tokens is the wrong denominator, and this book uses it only where nothing better exists. The economically meaningful quantity is the cost of a \emph{completed} unit of work:
\[
C_{\text{solved}} \;=\; \frac{C_{\text{input}} + C_{\text{output}} + C_{\text{tools}} + C_{\text{retries}}}{p_{\text{success}}},
\]
where parallel sampling counts every sampled token, and where wall-clock latency and maximum concurrency are reported alongside. The correction is not pedantic: reasoning-model output lengths have been inflating roughly fivefold per year~\cite{epoch-output}, so a model advertised at a tenth the price per output token that spends ten times the reasoning tokens per solved task has delivered no saving at all, and a headline ``150 times cheaper'' can survive translation into delivered task economics or vanish entirely. Wherever this book quotes a price ratio, read it as an upper bound on the true economic advantage until $C_{\text{solved}}$ is published---which, as Appendix~\ref{app:cards} records, essentially no vendor yet does.

\subsection{The economics: LLMflation}

Layered on falling hardware requirements is the steepest price deflation of any industrial input in history: inference cost at constant capability falls roughly 10$\times$ per year---GPT-3-level capability fell from \$60 to \$0.06 per million tokens in three years~\cite{llmflation}. Self-hosting break-evens are now measured in months for any serious deployment (roughly one month of payback at 50B tokens/month)~\cite{selfhost}. The strategic geometry is the solar module's: a Western closed-model business must recover \$100B-scale capex from a price that Chinese open weights plus community quantization reset toward marginal cost every quarter. What that geometry does not predict---and what the next section documents---is which party arrives at the low price first. Chapter~\ref{ch:trade} prices the consequences.

\section{The price inversion}
\label{sec:price-inversion}

The obvious next step in that argument is the one the evidence does not take. If open weights reset price toward marginal cost, the expected consequence is a Western frontier laboratory forced downward---and in July and August 2026 a Western frontier laboratory did move downward, sharply and in public. But the details of the move, and the direction the Chinese price leader moved at the same time, break the simple version of the story, and this section prints the break because the book's discipline is to print the datum that cuts against it.

Start with what was actually cut. On 30~July 2026---the press cycle ran on the 31st, which is why the date is often given a day late---OpenAI reduced the price of \textbf{GPT-5.6 Luna}, the cheapest tier of the family and the successor to the ``nano'' class, by 80\%: from \$1.00/\$6.00 to \$0.20/\$1.20 per million input/output tokens. Terra was cut 20\%, to \$2/\$12. \emph{Sol, the flagship, was not cut at all}, holding at \$5/\$30~\cite{openai-luna} [P]. Three weeks later, on 21~August, Sol went to \$4/\$20---explicitly ``for the next 3 months,'' which makes it a promotion rather than a repricing, and the distinction matters because a promotion is a defense of share while a repricing is a concession about cost. What the headline ``80\% price cut'' describes, then, is a deep cut at the bottom of the ladder, a shallow one in the middle, and nothing at the top: the shape of a firm defending the volume end of its business while leaving the capability rent intact.

\begin{table}[htbp]
  \centering
  \footnotesize
  \caption{The price inversion at \datafreeze: list price in US dollars per million tokens, ordered by the combined cost of one million input plus one million output tokens---a crude but consistent basis, and an upper bound on economic advantage rather than a cost per solved task (\S\ref{sec:cost-solved}). ``Open'' means downloadable weights, for which the listed price is the vendor's hosted endpoint and the marginal price to a self-hoster is hardware. $\dagger$~DeepSeek's July list, superseded by the increase of 13~August 2026~\cite{deepseek-raise}. $\ddagger$~Sol was \$5/\$30 until 21~August 2026---held unchanged through the 30~July cut---when it was cut to \$4/\$20 for a stated three months~\cite{openai-luna}.}
  \label{tab:price-inversion}
  \begin{tabular}{llcrrr}
    \toprule
    Model & Access & Origin & Input & Output & 1M+1M \\
    \midrule
    DeepSeek V4-Flash$^\dagger$ & open (MIT) & CN & 0.14 & 0.28 & 0.42 \\
    GLM-5.3-Flash & open (MIT) & CN & 0.15 & 0.50 & 0.65 \\
    MiniMax M2.5 & open (custom) & CN & 0.15 & 1.20 & 1.35 \\
    \textbf{GPT-5.6 Luna} & \textbf{closed API} & \textbf{US} & \textbf{0.20} & \textbf{1.20} & \textbf{1.40} \\
    GLM-5.3 & open (custom) & CN & 1.40 & 4.40 & 5.80 \\
    Qwen3.8-Max & open (custom) & CN & 2.00 & 6.00 & 8.00 \\
    GPT-5.6 Terra & closed API & US & 2.00 & 12.00 & 14.00 \\
    Kimi K3 & open (custom) & CN & 3.00 & 15.00 & 18.00 \\
    GPT-5.6 Sol$^\ddagger$ & closed API & US & 4.00 & 20.00 & 24.00 \\
    \bottomrule
  \end{tabular}
\end{table}

The consequence is the fact this chapter most needed and the one most easily missed. At \$0.20/\$1.20, GPT-5.6 Luna is \emph{cheaper on list price than most Chinese open-weight models}: below GLM-5.3, below Kimi K3, below Qwen3.8-Max, and below DeepSeek V4-Pro on input. On the combined million-in, million-out basis of Table~\ref{tab:price-inversion}, Luna costs \$1.40 against GLM-5.3's \$5.80, Kimi K3's \$18 and Claude Opus 5's \$30~\cite{price-inversion} [P/A]. The incumbent did not merely match the commodity price; at the bottom of its ladder it went underneath it. Two precisions keep that from becoming a slogan. The inversion holds against the \emph{frontier-scale} open models and not against the distilled tiers---GLM-5.3-Flash at \$0.65 combined and DeepSeek V4-Flash at its July list of \$0.42 both sit below Luna, and both are downloadable, so their marginal price to an organization with the hardware of \S\ref{sec:quant} is not a list price at all. And per-token list prices remain an upper bound on economic advantage rather than a measure of delivered work: a cheaper token that buys more tokens per solved task has bought nothing (\S\ref{sec:cost-solved}).

The cut worked, on the only evidence that can show it. Between 27~July and 14~August, Luna's usage rose to $13.8\times$ its pre-cut average and Terra's to $5.6\times$, against $1.1\times$ for the undiscounted Sol, with roughly 75\% of the gain taken from competitors rather than cannibalized from the vendor's own higher tiers~\cite{price-inversion} [A]. That is behavioural evidence of genuine cross-price elasticity at the volume end of the market---buyers moved, and they moved from somewhere else. The router data corroborate the shape: in the week to 29~August the fastest-growing Western entries on OpenRouter were its two cheapest metered tiers---GPT-5.6 Luna at 6.75 trillion tokens, up 31\% after the cut, and Gemini 3.7 Flash at 4.05 trillion, up 156\%---on a board still led by a Chinese open model, Ox~Alpha at 20.7 trillion and up 216\%, with DeepSeek V4 Flash at 12.3T and Xiaomi's MiMo-V2.5 at 10.0T behind it~\cite{openrouter-aug26} [A]. Router volumes are a small and promotion-distorted slice of global inference and are read here as a direction, not a census. One entry on that board belongs to neither camp as this chapter has drawn them: the sixth-largest model by routed tokens, at 5.62 trillion and up 27\%, is \emph{NVIDIA's} Nemotron 3 Ultra, released free and open-weight by the chip vendor whose accelerators serve most of the rest of the list. Chapter~\ref{ch:mirror} takes up what it means that the hardware incumbent is now the largest Western open-weight supplier.

Now the datum that cuts against the thesis. On 13~August 2026 DeepSeek---the laboratory that started the price war, the reference point for every ``150 times cheaper'' comparison in this literature, and the firm whose pricing this book has treated as the commodity floor---\emph{raised} its API prices by as much as 1{,}100\%, launching V4-Pro with peak and off-peak tiering; Caixin called it ``a break from the prolonged price war among China's AI developers''~\cite{deepseek-raise} [P]. No rationale was stated. The sequencing has to be read carefully, and it forbids the reading that would be most convenient in both directions: the deepest Western cut came on 30~July, two weeks \emph{before} the Chinese increase, so the cut cannot have been a response to it and the increase cannot have been a response to the cut. Candidate explanations exist---capacity scarcity making peak-hour tokens genuinely more expensive to serve, a deliberate move up-market as the distilled tiers take the volume, or a research house declining to subsidize inference it would rather spend on training---but each is inference, none is stated, and this book selects among them only to the extent of grading them all [M] and noting that the tiering itself is a capacity signal rather than a margin signal. What does not survive is the simple story in which cheap Chinese models force Western prices down: the price leader reversed direction mid-cascade, while the Western incumbent went below it. Commoditization at the volume end is real and observable in Table~\ref{tab:price-inversion}; the arrow of causation implied by the popular version of the story is not.

What the inversion does \emph{not} do is recover the commons, and the reason is the structural point of \S\ref{sec:maxtier} restated in price terms. A metered tier competes with a downloadable one only for the buyers who intend to rent tokens; for everyone who intends to run the model---the self-hosters of \S\ref{sec:quant}, the fine-tuners, the sovereign deployers of Chapter~\ref{ch:nonaligned}, the enterprises whose objection to a Chinese endpoint is jurisdictional rather than financial---the relevant price of an open model is not its vendor's list price at all, and no cut to a closed API changes it. Undercutting GLM-5.3's hosted endpoint therefore wins tokens without winning weights: it defends the volume tier of the rental business while leaving untouched the derivative ecosystems, the fine-tuning tooling, the quantization conventions and the standards gravity that Table~\ref{tab:metric-matrix} measures and that this chapter has argued are the currency the contest is actually denominated in. The two readings of the cut are correspondingly uncomfortable in opposite directions, and this book takes neither on faith: if \$0.20/\$1.20 is comfortably above cost, then the rent at the volume end was thin long before any Chinese release and the deflation argument of Chapter~\ref{ch:trade} must rest on the capability tier instead; if it is not, the incumbent is buying share at precisely the tier a commons with a floor at zero is best equipped to take. The evidence available at \datafreeze does not discriminate between them, which is why the dashboard watches the sequence of cuts rather than any instance.

The attribution discipline this book applies to Chinese vendor claims has to be applied symmetrically here, and it is uncomfortable. \emph{No OpenAI or Anthropic executive has attributed a 2026 price cut to Chinese competition on the record} [V]. Every company statement cites efficiency---including the striking and entirely unaudited claim, in OpenAI's own announcement of the 30~July cut, that GPT-5.6 Sol autonomously optimized its production kernels, ``reducing serving costs by 20\% and improving token-generation efficiency by 15\%''~\cite{openai-luna} [P, vendor]. The China attribution that appears throughout the trade press, and that this chapter's own argument would be pleased to accept, is press inference. The strongest adjacent primary statement is Altman's, and its provenance has to be stated, because it is not the pricing announcement: both remarks come from a podcast appearance on \emph{Invest Like the Best} on 29~July, the day before the cut. He said ``We better be the greatest and the cheapest,'' and, of Kimi~K3 specifically, ``You get a better deal today, at least at a particular like latency, using OpenAI's models than Kimi''---and the middle clause is load-bearing, because it restricts the comparison to a latency band rather than asserting a general price advantage. Read it precisely: that is a public commitment to undercut the Chinese open models, which is strong evidence that they are the competitive reference point---and it is not a concession that they caused the cut. The distinction is the difference between a firm that is watching a rival and a firm that has been forced by one, and the record supports the first without establishing the second.

Nor can this chapter claim the cuts are below cost, because nobody has shown it. The available unit economics are all reported via \emph{The Information} and unconfirmed by the companies: OpenAI at \$5.7B of first-quarter 2026 revenue against \$3.7B of cash burn, with adjusted gross margin around 33\%---though the reporting conflicts even on direction, one series having the figure fall from 46\% to 33\% and another rise from 33\% to 39\%, which is itself the point; Anthropic at a \$65B run rate at the end of July with \emph{API gross margin reported above 80\%}~\cite{openai-econ} [A]. \emph{No credible published estimate of loss per token exists, and no analyst asserts that these prices are below cost}---``inference loses money'' is a claim about model mix, free tiers and training amortization, not a universal one, and the 80\% figure sitting next to the 33\% figure shows how much of the disagreement is really about what is being counted. The honest statement of what the price line shows is therefore narrow: at the volume end, metered pricing has converged toward a level a commons with a floor at zero makes indefensible, and the American incumbent reached that level first and voluntarily; at the capability end, the rent is intact, unmatched, and---on the one margin figure anyone has published---very large. That is Proposition~P3 confirmed at the bottom of the ladder and unproven at the top, which is exactly where Chapter~\ref{ch:trade} takes it up.

\section{Openness as declared national strategy}

What began as firm-level catch-up strategy is now explicit state doctrine. The Global AI Governance Action Plan (WAIC, July 2025) commits China to ``cross-border open-source communities,'' open sharing of ``basic resources,'' and lowering ``the thresholds of technological innovation''---accompanied by the World AI Cooperation Organization, headquartered in Shanghai and founded in July 2026 with 29 member states (none a major Western democracy), a UN capacity-building group co-chaired with Zambia spanning 80 countries, and BRICS AI centers~\cite{action-plan,carnegie-pivot,waico}. The State Council's ``AI+'' plan targets self-controllable base models and a thousand deployed industry cases; draft policy counts open-source contribution toward academic credit~\cite{ai-plus}. As Tsinghua's Liu Zhiyuan put it: ``Open source has become politically correct''~\cite{mittr-open}. The export product is not a model; it is a \emph{stack}---weights, tooling, quantizations, and cheap tokens---aimed at the Global South, where two-thirds of EV sales already run on Chinese LFP and the AI adoption default is being set the same way. Even Western institutions that ban the DeepSeek \emph{app} on data-sovereignty grounds turn to self-hosted Chinese weights: Microsoft evaluated a fine-tuned DeepSeek V4 for a Copilot agent in June 2026~\cite{atlantic-council}.

\section{Compute constraints as forcing function}

Export controls shaped this element rather than preventing it. Denied H100-class bandwidth and capacity, Chinese labs innovated in exactly the directions that reduce dependence on both: extreme MoE sparsity (K3 activates 1.8\% of weights per token; V4-Pro 3\%), attention compression (MLA, DSA, KV caches at $\sim$2\% of dense), low-precision training (FP8, INT4, MXFP4), and multi-token prediction~\cite{raschka,kimi-k3}. The 2025 episode in which DeepSeek's next flagship was delayed by failed training runs on Huawei Ascend---then reverted to NVIDIA hardware while Ascend took inference---shows the transition is real but incomplete~\cite{ascend-delay}, and it is the reason this book grades the two halves of that transition separately. The report, five months later, that Zhipu had trained the GLM-5 family entirely on Ascend would show how fast the \emph{training} half is completing if it had a primary behind it; it does not---the claim circulates in the trade press without a source, the repository carries no hardware statement, and Huawei's own announcement describes post-release inference adaptation---so this book grades it [A/C], and the one Zhipu training claim that does have a primary, GLM-Image, discloses no scale and has not been audited~\cite{glm5,glm5-ascend-unsourced,zhipu-ascend-prior}. What is on the record by \datafreeze is the other half---a model that was globally top-four in its launch week \emph{served} at production scale on a domestic cluster, with no chipmaker named and no training claim attached (\S\ref{sec:glm53})---and the distinction is where the whole element turns, because inference is where the deployed volume, the standards and the revenue sit, while training is where the export controls actually bite. The design direction is now locked in: Chinese frontier models are co-evolving with the domestic hardware of Chapter~\ref{ch:compute}---sparse, capacity-hungry, bandwidth-frugal, natively low-precision---which is precisely the profile that CPU-class and LPDDR-based systems serve best. DeepSeek's adoption of the UE8M0 FP8 scale format ``for next-generation domestic chips'' made the co-design explicit~\cite{ue8m0}.

The model layer, in playbook terms, is already in the late-war stage, though the ledger has to be read line by line rather than declared. Adoption: captured. Standards---weight formats, sparsity recipes, RL pipelines, quantization conventions, and now the reference workloads against which American accelerators are benchmarked in public---increasingly set in Chinese repositories. Rent: destroyed at the volume end of the market, where Section~\ref{sec:price-inversion} finds list prices converged and the Western incumbent priced underneath the Chinese frontier, but plainly not at the capability end, where the only published margin figure in the record is above 80\%. What open weights cannot yet guarantee is the hardware to run and train them sovereignly---and this element closes with the serving half of that guarantee demonstrated at production scale and the training half resting on a press claim without a primary. That is the work of the next two elements.
